\section{Заключение}

Метод медленно меняющихся амплитуд позволяет заменить исследование слабонелинейной автоколебательной системы анализом укороченных уравнений. Состояния равновесия первого укороченного уравнения определяют возможные амплитуды автоколебаний и их устойчивость, а второе уравнение задаёт поправку к частоте.

Для лампового генератора характер режима возбуждения определяется положением рабочей точки и видом нечётной части нелинейной характеристики. При кубической аппроксимации реализуется мягкий режим, при пятой степени возможно возникновение жёсткого режима с гистерезисом, а при седьмой степени -- сложно-жёсткий режим с несколькими устойчивыми предельными циклами.

Экспериментальное исследование фазовой плоскости позволяет непосредственно наблюдать предельные циклы, спирали затухания и переходы между режимами, сопоставляя динамическую картину с результатами теоретического анализа.
